\section{Adversarial finite-state checks}
The examples below are exact toy models. They test universal statements and
assumptions; none is asserted to be the BLADE-NCO implementation.

\subsection{Stale memory after a target shift}
\label{sec:counterexamples}
\status{DISPROVED} for the claim that historical memory cannot hurt.
Take $\mathcal{X}=\{a,b\}$. Before a change the hard region is
$\HARD_0=\{a\}$; just before the first measured proposal it shifts to
$\HARD_i=\{b\}$ for all $i\geq1$. The inherited population contains only
$a$. The memory policy mutates from $a$ to $a$ with probability $0.99$ and to
$b$ with probability $0.01$, retaining $a$ after a miss. Each proposal costs
one unit. Its hitting time is geometric with success probability $0.01$,
so $\Ex[\tau_M]=100$ and
\[
 \Prb(\tau_M\leq B)=1-0.99^B.
\]
A fixed generator that independently samples $a$ or $b$ uniformly every
proposal has $\Ex[\tau_G]=2$ and success probability $1-0.5^B$. Thus the
memory policy is strictly worse at every positive integer budget. The target
has undergone a single non-stationary shift; the failure is stale structural
support, not an independence artifact.

\subsection{Gradient adaptation can beat structural mutation}
\status{DISPROVED} for any universal claim that structural mutation
outperforms HAC-style gradient adaptation. Let
$\mathcal{X}=\{0,1,\ldots,L\}$ and $\HARD=\{L\}$. A structural mutation
from $i<L$ advances to $i+1$ with probability $\varepsilon$ and otherwise
stays at $i$. Starting at $0$, its expected number of unit-cost proposals to
hit $L$ is $L/\varepsilon$. A gradient-informed comparator that advances
deterministically one step per unit-cost proposal hits in $L$ proposals.
For any $0<\varepsilon<1$, the structural policy is slower. This is a toy
landscape with an informative gradient, not a model of the actual HAC
implementation; it proves that operator labels alone imply no ranking.

\subsection{Diversity can allocate effort away from the target}
\status{DISPROVED} for the claim that increased diversity always improves
first hitting time. Let $\mathcal{X}=\{s,u,g\}$, target $\HARD=\{g\}$,
and one proposal cost one unit. The diverse parent set contains distinct
niches $s$ and $u$; parent selection is uniform. Mutation from $s$ reaches
$g$ surely, whereas mutation from $u$ stays at $u$. Hence the per-proposal
hit probability is $1/2$ and expected hitting cost is $2$. A focused policy
that selects $s$ every time hits in cost $1$. Preserving the irrelevant
niche does not help this target and reduces access to the useful parent.

\subsection{Surrogate error can remove a useful bridge}
\status{DISPROVED} for the claim that surrogate-based elite retention
preserves true-hardness reachability. Consider states $\{a,b,g\}$ and
current target $\HARD=\{g\}$; neither $a$ nor $b$ is itself a hit. From $a$,
mutation reaches $g$ surely; from $b$, mutation stays at $b$ surely. A
one-slot score-elite update retains whichever of $a,b$ has the larger
surrogate score. If the true scores of $a,b$ are equal but HGB40
systematically overestimates $b$, the update retains $b$, discards $a$, and
makes the target unreachable. An accurate or reachability-aware selection
would retain $a$. This is a counterexample to a universal error-robustness
claim, not a claim about the precise Elite-QD1000 update rule.

\subsection{Niche aggregation can destroy the Markov property}
\status{DISPROVED} for the assertion that a niche label alone is a Markov
state. Let full states $x_1,x_2$ both map to niche $A$, while $z,w$ map to
distinct niches $Z,W$. From $x_1$, the next state is $z$ with probability
$1$; from $x_2$, the next state is $w$ with probability $1$. Both full states
have the same current aggregate $A$, but the next aggregate is respectively
$Z$ or $W$. The transition probability out of aggregate $A$ is therefore
not determined by the aggregate alone (strong lumpability fails). A Markov
proof using only niche counts would be invalid unless the update rule supplies
the needed lumpability condition or the state is augmented with the hidden
structural/solver information.